\documentclass[20pt,a1paper,portrait]{tikzposter}
\usepackage[utf8]{inputenc}
\usepackage[T1]{fontenc}
\usepackage{lmodern}
\usepackage[spanish,es-nodecimaldot]{babel}
\usepackage{booktabs}
\usepackage{qrcode}
\usepackage{tikz}
\usetikzlibrary{arrows.meta,positioning,shapes.geometric}

\definecolor{azul}{HTML}{1F4E79}
\definecolor{claro}{HTML}{E8F0F8}
\definecolor{verde}{HTML}{2E8B57}
\definecolor{naranja}{HTML}{E08A1E}
\definecolor{gris}{HTML}{8A8F98}

\usetheme{Default}
\definecolorstyle{pcolors}{}{
  \colorlet{colorOne}{azul}\colorlet{colorTwo}{azul}\colorlet{colorThree}{azul}}{
  \colorlet{backgroundcolor}{white}\colorlet{framecolor}{azul}
  \colorlet{titlefgcolor}{white}\colorlet{titlebgcolor}{azul}
  \colorlet{blocktitlebgcolor}{azul}\colorlet{blocktitlefgcolor}{white}
  \colorlet{blockbodybgcolor}{claro}\colorlet{innerblocktitlebgcolor}{azul}
  \colorlet{innerblocktitlefgcolor}{white}\colorlet{innerblockbodybgcolor}{claro}
  \colorlet{notefgcolor}{black}\colorlet{notebgcolor}{claro}\colorlet{noteframecolor}{azul}}
\usecolorstyle{pcolors}

\title{\textbf{Extractor de facturas FEL con LLM}}
\author{Angel Esteban Esquit Hernández y Javier Eduardo España Pacheco}
\institute{CC3092 Deep Learning -- Proyecto 3: Casos de uso de LLM \\ Feria científica, 2 de noviembre de 2026}

\begin{document}
\maketitle

\begin{columns}
\column{0.5}
\block{Problema}{
  Pymes y contadores en Guatemala reciben facturas FEL en PDF o imagen, con layouts distintos, y las digitan a mano en hojas de cálculo.
  Es lento y propenso a errores.\\[3mm]
  \textbf{Objetivo:} subir una factura y obtener un JSON/CSV revisable, con los campos dudosos marcados para revisión humana.\\[3mm]
  \textbf{¿Por qué un LLM?} Las reglas (regex) no escalan a layouts, ruido y fotos: el baseline solo acierta el 23\,\% de los campos.
}

\block{Pipeline}{
  \begin{center}
  \begin{tikzpicture}[
    box/.style={draw=azul, fill=white, rounded corners, align=center, minimum height=20mm, minimum width=50mm, thick},
    arrow/.style={-{Latex[length=4mm]}, very thick, azul}]
    \node[box] (a) {PDF o\\imagen};
    \node[box, right=14mm of a] (b) {LLM\\salida estructurada\\(esquema Pydantic)};
    \node[box, below=14mm of b] (c) {Validadores\\deterministas\\IVA, sumas, NIT, fecha, UUID};
    \node[box, left=14mm of c] (d) {Interfaz Streamlit\\tabla, alertas,\\JSON / CSV};
    \draw[arrow] (a) -- (b);
    \draw[arrow] (b) -- (c);
    \draw[arrow] (c) -- (d);
    \draw[arrow, naranja] (c.east) -- ++(10mm,0) |- node[pos=0.75, above, align=center, font=\small, text=black]{reintento (máx. 1)} (b.east);
  \end{tikzpicture}
  \end{center}
  Modelos comparados: \texttt{gpt-4o-mini} (OpenAI), \texttt{claude-haiku-5-5} y \texttt{claude-sonnet-5-5} (Anthropic).
  Técnica: prompting + esquema + validación; sin RAG ni fine-tuning.
}

\block{Datos de evaluación}{
  \begin{itemize}
    \item \textbf{45 facturas FEL sintéticas} (semilla 2026): 4 layouts, PDF e imagen, 3 formatos de fecha, 1--8 líneas, ruido, rotación y desenfoque.
    \item \textbf{15 fixtures públicas} (Shopify Invoice Test Pack, CC BY 4.0): multipágina, árabe, descuentos, flete, mala calidad.
  \end{itemize}
  Prompt v1 mínimo $\rightarrow$ v2 con reglas de nulos, normalización y defensa contra instrucciones incrustadas en el documento.
}

\block{Pruébalo}{
  \begin{minipage}{0.28\linewidth}
    \centering
    \qrcode[height=5cm]{https://github.com/Javier-Espana/PRY3-DL}
  \end{minipage}\hfill
  \begin{minipage}{0.68\linewidth}
    \textbf{Código, README y reproducción de métricas:}\\
    github.com/Javier-Espana/PRY3-DL
  \end{minipage}
}
\column{0.5}
\block{Resultados principales: 45 facturas FEL}{
  \begin{center}
  \begin{tikzpicture}[x=1mm,y=1mm]
    \def\w{1.3}
    \foreach \n/\v/\c/\i in {Regex/23.25/gris/0, GPT-4o-mini/90.77/naranja/1, Haiku 5.5/100/verde/2, Sonnet 5.5/98.63/azul/3}{
      \pgfmathsetmacro{\y}{-\i*17}
      \node[anchor=east, align=right] at (-2,\y) {\n};
      \fill[\c] (0,\y-5) rectangle (\v*\w,\y+5);
      \node[anchor=west] at (\v*\w+2,\y) {\v\,\%};
    }
    \draw[thick] (0,6) -- (0,-56);
  \end{tikzpicture}\\[1mm]
  {\small Exactitud por campo, sin reintentos}
  \end{center}
  \begin{center}
  \begin{tabular}{lrrr}
    \toprule
    Sistema & Docs. 100\,\% & Latencia & USD/doc. \\
    \midrule
    Regex & 0\,\% & -- & 0 \\
    GPT-4o-mini & 15,56\,\% & 4,9 s & 0,00147 \\
    Haiku 5.5 & \textbf{100\,\%} & 5,3 s & \textbf{0,00076} \\
    Sonnet 5.5 & 82,22\,\% & 4,0 s & 0,0113 \\
    \bottomrule
  \end{tabular}
  \end{center}
  \textbf{En el conjunto FEL, Haiku 5.5 fue la mejor relación precisión/costo}: 100\,\% de exactitud a USD 0,00076 por documento. Sonnet costó unas 15 veces más y obtuvo 98,63\,\%. Costo total de las tres corridas: $\approx$ USD 0,61 (prompt v2, 0 reintentos, 2.000 tokens de salida).
}

\block{Robustez y hallazgos}{
  \begin{itemize}
    \item \textbf{Prueba secundaria (15 fixtures públicas internacionales):} GPT-4o-mini 72,31\,\%, Haiku 83,08\,\% y Sonnet 91,28\,\% (USD 0,01162 vs. 0,00078 por documento). Sonnet fue mejor allí, a ~15 veces el costo. Son documentos árabes, bilingües, con descuentos y flete; no miden desempeño FEL.
    \item \textbf{La elección del modelo depende del dominio y del presupuesto:} la recomendación de Haiku es específica del conjunto FEL evaluado.
    \item \textbf{Error de etiquetado hallado y corregido:} 6 facturas sin IVA tenían \texttt{subtotal} en el gold sin que el documento lo mostrara.
    \item \textbf{Compromiso del prompt v2:} ``null si no aparece'' evita inventar, pero deja vacía la moneda cuando solo dice ``Q''.
  \end{itemize}
}

\block{Limitaciones}{
  \begin{itemize}
    \item El conjunto principal es sintético y lo creó el mismo equipo que el extractor: techo de rendimiento probable.
    \item Falta validar con facturas reales anonimizadas y con usuarios externos.
    \item Las facturas se envían al proveedor del modelo: no subir datos que no puedan compartirse.
    \item Los validadores detectan inconsistencias, no garantizan exactitud fiscal; siempre hay revisión humana.
  \end{itemize}
}

\end{columns}

\end{document}
